\section{Objetivos del proyecto}

\subsection{Objetivo general}
El objetivo general de este proyecto consiste en diseñar, desarrollar y evaluar computacionalmente un sistema de Inteligencia Artificial capaz de clasificar patrones vocales asociados a enfermedades neurodegenerativas. El propósito es generar una herramienta tecnológica de bajo coste y alta accesibilidad que funcione como sistema de apoyo a la decisión clínica (CDSS), permitiendo identificar biomarcadores acústicos subclínicos en etapas tempranas de la enfermedad.

\subsection{Objetivos específicos}
Para garantizar la consecución del objetivo general, el proyecto se estructura en los siguientes objetivos específicos, definidos bajo criterios de viabilidad técnica y rigor metodológico:

\begin{itemize}
    \item \textbf{Analizar el estado del arte:} Realizar una revisión sistemática de la literatura científica sobre el uso de la voz como biomarcador clínico y estudiar las arquitecturas computacionales actuales empleadas en este dominio.
    
    \item \textbf{Preprocesar y extraer características acústicas:} Diseñar un pipeline de procesamiento de señales de audio a partir de bases de datos clínicas estandarizadas (NeuroVoz y PC-GITA), generando tanto características acústicas manuales para modelos clásicos como representaciones espectrales (espectrogramas) para redes neuronales profundas.
    
    \item \textbf{Desarrollar y comparar arquitecturas de Inteligencia Artificial:} Implementar y entrenar distintos enfoques de modelado, evaluando comparativamente el rendimiento de algoritmos de Machine Learning tradicional (KNN, XGBoost), arquitecturas de Deep Learning basadas en visión (ResNet18) y enfoques de vanguardia basados en Foundation Models (Wav2Vec 2.0).
    
    \item \textbf{Evaluar el rendimiento y la interpretabilidad clínica:} Cuantificar la eficacia de los modelos mediante métricas de evaluación estándar (precisión, sensibilidad, especificidad) e integrar técnicas de Inteligencia Artificial Explicable (como valores SHAP) para garantizar que las predicciones sean transparentes e interpretables por el personal médico.
    
    \item \textbf{Construir una prueba de concepto interactiva y desplegable:} Desarrollar una aplicación cliente-servidor (frontend móvil/web y API REST) que integre los modelos entrenados, proporcionando una interfaz intuitiva para la evaluación de pacientes en tiempo real.
    
    \item \textbf{Garantizar la reproducibilidad del proyecto:} Documentar, empaquetar y liberar tanto el código fuente como los binarios generados bajo un entorno de control de versiones público, fomentando la transparencia y la continuidad investigadora.
\end{itemize}